% Options for packages loaded elsewhere
\PassOptionsToPackage{unicode}{hyperref}
\PassOptionsToPackage{hyphens}{url}
\documentclass[
  ignorenonframetext,
]{beamer}
\newif\ifbibliography
\usepackage{pgfpages}
\setbeamertemplate{caption}[numbered]
\setbeamertemplate{caption label separator}{: }
\setbeamercolor{caption name}{fg=normal text.fg}
\beamertemplatenavigationsymbolsempty
% remove section numbering
\setbeamertemplate{part page}{
  \centering
  \begin{beamercolorbox}[sep=16pt,center]{part title}
    \usebeamerfont{part title}\insertpart\par
  \end{beamercolorbox}
}
\setbeamertemplate{section page}{
  \centering
  \begin{beamercolorbox}[sep=12pt,center]{section title}
    \usebeamerfont{section title}\insertsection\par
  \end{beamercolorbox}
}
\setbeamertemplate{subsection page}{
  \centering
  \begin{beamercolorbox}[sep=8pt,center]{subsection title}
    \usebeamerfont{subsection title}\insertsubsection\par
  \end{beamercolorbox}
}
% Prevent slide breaks in the middle of a paragraph
\widowpenalties 1 10000
\raggedbottom
\AtBeginPart{
  \frame{\partpage}
}
\AtBeginSection{
  \ifbibliography
  \else
    \frame{\sectionpage}
  \fi
}
\AtBeginSubsection{
  \frame{\subsectionpage}
}
\usepackage{iftex}
\ifPDFTeX
  \usepackage[T1]{fontenc}
  \usepackage[utf8]{inputenc}
  \usepackage{textcomp} % provide euro and other symbols
\else % if luatex or xetex
  \usepackage{unicode-math} % this also loads fontspec
  \defaultfontfeatures{Scale=MatchLowercase}
  \defaultfontfeatures[\rmfamily]{Ligatures=TeX,Scale=1}
\fi
\usepackage{lmodern}
\ifPDFTeX\else
  % xetex/luatex font selection
\fi
% Use upquote if available, for straight quotes in verbatim environments
\IfFileExists{upquote.sty}{\usepackage{upquote}}{}
\IfFileExists{microtype.sty}{% use microtype if available
  \usepackage[]{microtype}
  \UseMicrotypeSet[protrusion]{basicmath} % disable protrusion for tt fonts
}{}
\makeatletter
\@ifundefined{KOMAClassName}{% if non-KOMA class
  \IfFileExists{parskip.sty}{%
    \usepackage{parskip}
  }{% else
    \setlength{\parindent}{0pt}
    \setlength{\parskip}{6pt plus 2pt minus 1pt}}
}{% if KOMA class
  \KOMAoptions{parskip=half}}
\makeatother
\usepackage{color}
\usepackage{fancyvrb}
\newcommand{\VerbBar}{|}
\newcommand{\VERB}{\Verb[commandchars=\\\{\}]}
\DefineVerbatimEnvironment{Highlighting}{Verbatim}{commandchars=\\\{\}}
% Add ',fontsize=\small' for more characters per line
\newenvironment{Shaded}{}{}
\newcommand{\AlertTok}[1]{\textcolor[rgb]{1.00,0.00,0.00}{\textbf{#1}}}
\newcommand{\AnnotationTok}[1]{\textcolor[rgb]{0.38,0.63,0.69}{\textbf{\textit{#1}}}}
\newcommand{\AttributeTok}[1]{\textcolor[rgb]{0.49,0.56,0.16}{#1}}
\newcommand{\BaseNTok}[1]{\textcolor[rgb]{0.25,0.63,0.44}{#1}}
\newcommand{\BuiltInTok}[1]{\textcolor[rgb]{0.00,0.50,0.00}{#1}}
\newcommand{\CharTok}[1]{\textcolor[rgb]{0.25,0.44,0.63}{#1}}
\newcommand{\CommentTok}[1]{\textcolor[rgb]{0.38,0.63,0.69}{\textit{#1}}}
\newcommand{\CommentVarTok}[1]{\textcolor[rgb]{0.38,0.63,0.69}{\textbf{\textit{#1}}}}
\newcommand{\ConstantTok}[1]{\textcolor[rgb]{0.53,0.00,0.00}{#1}}
\newcommand{\ControlFlowTok}[1]{\textcolor[rgb]{0.00,0.44,0.13}{\textbf{#1}}}
\newcommand{\DataTypeTok}[1]{\textcolor[rgb]{0.56,0.13,0.00}{#1}}
\newcommand{\DecValTok}[1]{\textcolor[rgb]{0.25,0.63,0.44}{#1}}
\newcommand{\DocumentationTok}[1]{\textcolor[rgb]{0.73,0.13,0.13}{\textit{#1}}}
\newcommand{\ErrorTok}[1]{\textcolor[rgb]{1.00,0.00,0.00}{\textbf{#1}}}
\newcommand{\ExtensionTok}[1]{#1}
\newcommand{\FloatTok}[1]{\textcolor[rgb]{0.25,0.63,0.44}{#1}}
\newcommand{\FunctionTok}[1]{\textcolor[rgb]{0.02,0.16,0.49}{#1}}
\newcommand{\ImportTok}[1]{\textcolor[rgb]{0.00,0.50,0.00}{\textbf{#1}}}
\newcommand{\InformationTok}[1]{\textcolor[rgb]{0.38,0.63,0.69}{\textbf{\textit{#1}}}}
\newcommand{\KeywordTok}[1]{\textcolor[rgb]{0.00,0.44,0.13}{\textbf{#1}}}
\newcommand{\NormalTok}[1]{#1}
\newcommand{\OperatorTok}[1]{\textcolor[rgb]{0.40,0.40,0.40}{#1}}
\newcommand{\OtherTok}[1]{\textcolor[rgb]{0.00,0.44,0.13}{#1}}
\newcommand{\PreprocessorTok}[1]{\textcolor[rgb]{0.74,0.48,0.00}{#1}}
\newcommand{\RegionMarkerTok}[1]{#1}
\newcommand{\SpecialCharTok}[1]{\textcolor[rgb]{0.25,0.44,0.63}{#1}}
\newcommand{\SpecialStringTok}[1]{\textcolor[rgb]{0.73,0.40,0.53}{#1}}
\newcommand{\StringTok}[1]{\textcolor[rgb]{0.25,0.44,0.63}{#1}}
\newcommand{\VariableTok}[1]{\textcolor[rgb]{0.10,0.09,0.49}{#1}}
\newcommand{\VerbatimStringTok}[1]{\textcolor[rgb]{0.25,0.44,0.63}{#1}}
\newcommand{\WarningTok}[1]{\textcolor[rgb]{0.38,0.63,0.69}{\textbf{\textit{#1}}}}
\setlength{\emergencystretch}{3em} % prevent overfull lines
\providecommand{\tightlist}{%
  \setlength{\itemsep}{0pt}\setlength{\parskip}{0pt}}
% Slide layout defaults for warning-clean scientific decks.
\usepackage{etoolbox}
\IfFileExists{xurl.sty}{\usepackage{xurl}}{}
\IfFileExists{seqsplit.sty}{\usepackage{seqsplit}}{\newcommand{\seqsplit}[1]{#1}}
\protected\def\breakseq#1{\seqsplit{#1}}
\protected\def\breaktt#1{\begingroup\ttfamily\seqsplit{#1}\endgroup}
\setlength{\emergencystretch}{6em}
\tolerance=5000
\hbadness=10000
\hfuzz=1pt
\setlength{\tabcolsep}{2pt}
\AtBeginEnvironment{longtable}{\tiny\renewcommand{\arraystretch}{0.86}\setlength{\tabcolsep}{1pt}}
\AtBeginEnvironment{tabular}{\tiny\renewcommand{\arraystretch}{0.86}\setlength{\tabcolsep}{1pt}}

% Natbib and cross-reference fallbacks — slides don't load natbib
% or cleveref, but combined-PDF manuscript prose may emit these
% commands. The fallback renders citations as a bracketed key list
% and cross-references as detokenized labels so slides don't fail on
% undefined control sequences or raw underscores. \providecommand is
% a no-op if packages load later.
\providecommand{\citep}[1]{[#1]}
\providecommand{\citet}[1]{#1}
\providecommand{\citealp}[1]{#1}
\providecommand{\citeauthor}[1]{#1}
\providecommand{\citeyear}[1]{#1}
\providecommand{\cref}[1]{\texttt{\detokenize{#1}}}
\providecommand{\Cref}[1]{\texttt{\detokenize{#1}}}

\newtheorem{proposition}{Proposition}
\newtheorem{hypothesis}{Hypothesis}
\usepackage{bookmark}
\IfFileExists{xurl.sty}{\usepackage{xurl}}{} % add URL line breaks if available
\urlstyle{same}
% fallback for those not using the hyperref driver hyperxmp:
\makeatletter
\@ifundefined{xmpquote}{\newcommand{\xmpquote}[1]{#1}}{}
\makeatother
\hypersetup{
  hidelinks,
  pdfcreator={LaTeX via pandoc}}

\author{\texorpdfstring{}{}}
\date{}

\begin{document}

\section{\texorpdfstring{Appendix: Mathematical and Algorithmic Details
\label{sec:technical_appendix}}{Appendix: Mathematical and Algorithmic Details }}\label{appendix-mathematical-and-algorithmic-details}

\begin{frame}[allowframebreaks]{Appendix: Mathematical and Algorithmic
Details \label{sec:technical_appendix}}
\emph{This appendix collects the formal mathematical definitions,
derivations, and algorithmic specifications referenced from the main
methodology section. Each subsection is self-contained; equations are
labelled for cross-referencing from the body and from
\S\ref{tab:notation_symbols}.}
\end{frame}

\begin{frame}[fragile,allowframebreaks]{Citation-Weighted Hypothesis
Scoring Formula \label{sec:appendix_scoring}}
\protect\phantomsection\label{citation-weighted-hypothesis-scoring-formula}
For each hypothesis \(H\), we compute a citation-weighted evidence score
aggregating all assertions relevant to \(H\):

\begin{equation}
\score(H) = \frac{\sum_{a \in S(H)} w(a) - \sum_{a \in C(H)} w(a)}{\sum_{a \in A(H)} w(a)} \label{eq:app_score}
\end{equation}

where \(S(H)\) is the set of supporting assertions, \(C(H)\) the set of
contradicting assertions, \(A(H)\) all assertions for \(H\) (including
neutral), and the weight function is

\begin{equation}
w(a) = \log(1 + \text{citations}(a)) \cdot \text{confidence}(a). \label{eq:app_weight}
\end{equation}

The logarithmic citation weighting assigns higher weight to papers that
have accumulated more citations within the retrieved corpus window. This
measures \textbf{published attention and visibility} in the literature
graph---not truth, popularity on social media, field size alone, or
citation-network centrality (PageRank is not used). We report
sensitivity analyses under uniform weighting, confidence-only, raw
citation (popularity stress test), age discount, and field-normalized
cohort weighting; rank-stability Spearman \(\rho = 0.762\) versus the
default policy. The score lies in \([-1, 1]\): values indicate
citation-weighted evidence mapping and hypothesis triage within the
corpus, not calibrated confirmation of scientific truth.

\textbf{Temporal aggregation.} We additionally compute temporal trends
by evaluating the cumulative score at each year \(t\), using only
assertions from papers published in year \(\leq t\):

\begin{equation}
\score(H, t) = \frac{\sum_{a \in S(H,t)} w(a) - \sum_{a \in C(H,t)} w(a)}{\sum_{a \in A(H,t)} w(a)}. \label{eq:app_score_t}
\end{equation}

This reveals whether support for a hypothesis is growing, declining, or
plateauing over time. Cumulative aggregation (rather than yearly
windows) is preferred because per-year assertion counts for narrow
hypotheses are too sparse for stable point estimates.

\textbf{Algorithmic specification.} The scoring routine is a pure
function of the assertion set; it has no hidden state and is
deterministic given the input. The reference implementation lives in
\breaktt{src/knowledge\_graph/hypothesis.py} with weight policies in
\breaktt{src/knowledge\_graph/hypothesis\_weights.py}:

\begin{Shaded}
\begin{Highlighting}[]
\NormalTok{function score(H, assertions):}
\NormalTok{    S, C, A\_all = 0, 0, 0}
\NormalTok{    for a in assertions where a.hypothesis == H:}
\NormalTok{        w = log(1 + a.citations) * a.confidence}
\NormalTok{        if a.direction == "supports":     S     += w}
\NormalTok{        elif a.direction == "contradicts": C     += w}
\NormalTok{        A\_all += w}
\NormalTok{    return (S {-} C) / A\_all  if A\_all \textgreater{} 0  else 0}
\end{Highlighting}
\end{Shaded}

Boundary tests in \breaktt{tests/test\_scoring.py} confirm that
all-support fixtures yield \(+1\), all-contradict fixtures yield \(-1\),
and balanced fixtures yield \(0\) within numerical tolerance.

\FloatBarrier
\end{frame}

\begin{frame}[allowframebreaks]{Non-negative Matrix Factorization (NMF)
for Topic Modeling \label{sec:appendix_nmf}}
\protect\phantomsection\label{non-negative-matrix-factorization-nmf-for-topic-modeling}
We apply NMF to the TF-IDF matrix of the corpus to discover latent
topics. Given the document-term matrix
\(V \in \mathbb{R}^{n \times m}_{\geq 0}\), NMF finds factor matrices
\(W \in \mathbb{R}^{n \times k}_{\geq 0}\) and
\(H \in \mathbb{R}^{k \times m}_{\geq 0}\) such that \(V \approx W H\),
where \(k\) is the number of topics. We use multiplicative update rules
\citep{lee1999nmf}:

\begin{equation}
H \leftarrow H \odot \frac{W^T V}{W^T W H + \epsilon}, \qquad W \leftarrow W \odot \frac{V H^T}{W H H^T + \epsilon}, \label{eq:nmf_update}
\end{equation}

with \(\epsilon = 10^{-10}\) for numerical stability and a fixed random
seed of 42 for reproducibility (deterministic topic alignment across
pipeline runs, with empirical stability confirmed via Jaccard
similarities \(> 0.90\) across alternative seeds).

\textbf{Term-Frequency Inverse Document Frequency (TF-IDF).} The
document-term matrix is constructed using a smoothed TF-IDF weighting
\citep{salton1975vector}. For term \(t\) in document \(d\):

\begin{equation}
\text{TF-IDF}(t, d) = \text{tf}(t, d) \cdot \left[\log\!\left(\frac{N}{\text{df}(t) + 1}\right) + 1\right], \label{eq:tfidf}
\end{equation}

where \(\text{tf}(t, d) = \text{count}(t,d) / |d|\) is the normalized
term frequency, \(N\) the total number of documents, and
\(\text{df}(t)\) the document frequency of term \(t\). The \(+1\)
additive smoothing in the denominator prevents division by zero and
reduces the weight of extremely rare terms; the outer \(+1\) ensures
strictly positive IDF values. Document vectors are L2-normalized before
NMF factorization.

\FloatBarrier
\end{frame}

\begin{frame}[allowframebreaks]{Field Growth-Rate Estimation
\label{sec:appendix_growth}}
\protect\phantomsection\label{field-growth-rate-estimation}
The \textbf{mean year-over-year growth rate} \(\bar{g}\) is the
arithmetic mean of annual growth rates computed only for years where the
prior year had non-zero publications:

\begin{equation}
\bar{g} = \frac{1}{|Y|} \sum_{y \in Y} \frac{n_y - n_{y-1}}{n_{y-1}}, \label{eq:mean_growth}
\end{equation}

where \(Y = \{y : n_{y-1} > 0\}\) and \(n_y\) is the number of
publications in year \(y\).

The \textbf{doubling time} \(t_d\) is derived from the mean annual
growth rate:

\begin{equation}
t_d = \frac{\ln 2}{\ln(1 + \bar{g})}. \label{eq:doubling_time}
\end{equation}

The \textbf{compound annual growth rate} (CAGR) over the full span
\([y_0, y_T]\) is

\begin{equation}
\cagr = \left(\frac{n_{\text{cumulative}}(y_T)}{n_{\text{cumulative}}(y_0)}\right)^{1/(y_T - y_0)} - 1. \label{eq:cagr}
\end{equation}

For the current corpus, \(\cagr = 20.36\%\). The more recent growth
phase (2010--2026) exhibits substantially higher annualized growth than
the long-run average; reporting both the \(T\)-year CAGR and
recent-phase CAGR avoids overstating maturity-era expansion.

\FloatBarrier
\end{frame}

\begin{frame}[allowframebreaks]{Advanced Visualization Methods
\label{sec:appendix_viz}}
\protect\phantomsection\label{advanced-visualization-methods}
\begin{block}{PCA of TF-IDF Embeddings}
\protect\phantomsection\label{pca-of-tf-idf-embeddings}
Principal Component Analysis (PCA) is applied to the TF-IDF matrix \(V\)
to project each document into a 2-D space. The projection preserves the
directions of maximum variance, enabling visual inspection of document
clustering by domain. Loading arrows overlay the top-variance terms onto
the scatter plot, showing which vocabulary drives the principal
components.
\end{block}

\begin{block}{Hierarchical Clustering Dendrogram}
\protect\phantomsection\label{hierarchical-clustering-dendrogram}
For each domain \(s\), we compute the centroid
\(\bar{v}_s = \frac{1}{|D_s|} \sum_{d \in D_s} v_d\) where \(D_s\) is
the set of documents in domain \(s\) and \(v_d\) is the TF-IDF vector of
document \(d\). Ward linkage is applied to the centroid matrix to
produce a hierarchical clustering dendrogram showing semantic proximity
between domains.
\end{block}

\begin{block}{Term Heatmap}
\protect\phantomsection\label{term-heatmap}
For each domain \(s\) and term \(t\), we compute the mean TF-IDF weight
\(\bar{w}_{s,t} = \frac{1}{|D_s|} \sum_{d \in D_s} \text{TF-IDF}(t, d)\).
The heatmap displays \(\bar{w}_{s,t}\) for the top-\(k\) terms (by
global document frequency) across all domains, with cell intensity
proportional to mean weight. This reveals distinctive vocabulary
patterns that differentiate domains beyond the keyword-level
classification used for subfield assignment.
\end{block}

\begin{block}{Term Co-occurrence Matrix}
\protect\phantomsection\label{term-co-occurrence-matrix}
The co-occurrence matrix \(C \in \mathbb{R}^{k \times k}\) counts the
number of documents in which two terms appear together. For top-\(k\)
terms by document frequency,
\(C_{ij} = |\{d : t_i \in d \land t_j \in d\}|\). The matrix is
normalized to \([0, 1]\) by dividing by the maximum entry and visualized
as a symmetric heatmap.

\FloatBarrier
\end{block}
\end{frame}

\begin{frame}[fragile,allowframebreaks]{Nanopublication RDF Schema
\label{sec:appendix_rdf}}
\protect\phantomsection\label{nanopublication-rdf-schema}
Each nanopublication is serialized to RDF/TriG per the nanopublication
standard \citep{groth2010anatomy, kuhn2016decentralized}, producing four
named graphs. The following annotated example illustrates the structure
for a single assertion:

\begin{Shaded}
\begin{Highlighting}[]
\NormalTok{@prefix np:   \textless{}http://www.nanopub.org/nschema\#\textgreater{} .}
\NormalTok{@prefix prov: \textless{}http://www.w3.org/ns/prov\#\textgreater{} .}
\NormalTok{@prefix dc:   \textless{}http://purl.org/dc/terms/\textgreater{} .}
\NormalTok{@prefix aif:  \textless{}http://activeinference.institute/ontology/\textgreater{} .}
\NormalTok{@prefix xsd:  \textless{}http://www.w3.org/2001/XMLSchema\#\textgreater{} .}

\NormalTok{\# HEAD GRAPH: links nanopub to its three component graphs}
\NormalTok{\textless{}http://activeinference.institute/nanopub/a1b2c3d4e5f6\#head\textgreater{} \{}
\NormalTok{  \textless{}http://activeinference.institute/nanopub/a1b2c3d4e5f6\textgreater{}}
\NormalTok{    a np:Nanopublication ;}
\NormalTok{    np:hasAssertion       \textless{}...\#assertion\textgreater{} ;}
\NormalTok{    np:hasProvenance      \textless{}...\#provenance\textgreater{} ;}
\NormalTok{    np:hasPublicationInfo \textless{}...\#pubinfo\textgreater{} .}
\NormalTok{\}}

\NormalTok{\# ASSERTION GRAPH: the core scientific claim}
\NormalTok{\textless{}http://activeinference.institute/nanopub/a1b2c3d4e5f6\#assertion\textgreater{} \{}
\NormalTok{  aif:paper/10.1038\_nrn2787 aif:asserts aif:assertion/a1b2c3 .}
\NormalTok{  aif:assertion/a1b2c3}
\NormalTok{    aif:supports       aif:hypothesis/fep\_universality ;}
\NormalTok{    aif:claim          "The paper provides foundational support for FEP as a}
\NormalTok{                        unified brain theory."\^{}\^{}xsd:string ;}
\NormalTok{    aif:confidence     "0.85"\^{}\^{}xsd:double ;}
\NormalTok{    aif:citationCount  "5000"\^{}\^{}xsd:integer .}
\NormalTok{\}}

\NormalTok{\# PROVENANCE GRAPH: extraction lineage}
\NormalTok{\textless{}http://activeinference.institute/nanopub/a1b2c3d4e5f6\#provenance\textgreater{} \{}
\NormalTok{  aif:assertion/a1b2c3}
\NormalTok{    prov:wasGeneratedBy   \textless{}http://activeinference.institute/nanopub/a1b2c3d4e5f6\textgreater{} ;}
\NormalTok{    prov:generatedAtTime  "2026{-}01{-}15T12:00:00+00:00"\^{}\^{}xsd:dateTime ;}
\NormalTok{    prov:wasAttributedTo  "act\_inf\_metaanalysis/gemma3:4b"\^{}\^{}xsd:string ;}
\NormalTok{    prov:hadPrimarySource aif:paper/10.1038\_nrn2787 .}
\NormalTok{\}}

\NormalTok{\# PUBLICATION INFO GRAPH: nanopublication metadata}
\NormalTok{\textless{}http://activeinference.institute/nanopub/a1b2c3d4e5f6\#pubinfo\textgreater{} \{}
\NormalTok{  \textless{}http://activeinference.institute/nanopub/a1b2c3d4e5f6\textgreater{}}
\NormalTok{    dc:created "2026{-}01{-}15T12:00:00+00:00"\^{}\^{}xsd:dateTime ;}
\NormalTok{    dc:creator "act\_inf\_metaanalysis/gemma3:4b"\^{}\^{}xsd:string ;}
\NormalTok{    dc:license \textless{}https://creativecommons.org/publicdomain/zero/1.0/\textgreater{} .}
\NormalTok{\}}
\end{Highlighting}
\end{Shaded}

\begin{block}{Namespace Definitions}
\protect\phantomsection\label{namespace-definitions}
\begin{table}[H]
\centering
\caption{RDF namespace definitions used in the knowledge graph and nanopublication serialization. Each prefix maps to a W3C or domain-specific URI.}
\label{tab:namespace_definitions}
\begin{tabular}{lll}
\toprule
\textbf{Prefix} & \textbf{URI} & \textbf{Purpose} \\
\midrule
\texttt{np:}   & \breaktt{http://www.nanopub.org/nschema\#}     & Nanopub structural predicates \\
\texttt{prov:} & \breaktt{http://www.w3.org/ns/prov\#}          & PROV-O provenance model \\
\texttt{dc:}   & \breaktt{http://purl.org/dc/terms/}            & Dublin Core metadata \\
\texttt{aif:}  & \breaktt{http://activeinference.institute/ontology/} & Domain-specific predicates \\
\texttt{xsd:}  & \breaktt{http://www.w3.org/2001/XMLSchema\#}   & XML Schema datatypes \\
\bottomrule
\end{tabular}
\end{table}
\end{block}

\begin{block}{Core Triple Patterns}
\protect\phantomsection\label{core-triple-patterns}
The knowledge graph encodes five fundamental relationships:

\begin{table}[H]
\centering
\caption{Core RDF triple patterns encoding the five fundamental relationships in the knowledge graph. Each pattern links paper, assertion, hypothesis, or subfield nodes.}
\label{tab:core_triple_patterns}
\begin{tabular}{ll}
\toprule
\textbf{Triple Pattern} & \textbf{Meaning} \\
\midrule
\breaktt{Paper --aif:asserts--> Assertion}        & A paper makes a claim \\
\breaktt{Paper --aif:cites--> Paper}              & Intra-corpus citation link \\
\breaktt{Paper --aif:belongsTo--> Subfield}       & Domain classification \\
\breaktt{Assertion --aif:supports--> Hypothesis}  & Supporting evidence \\
\breaktt{Assertion --aif:contradicts--> Hypothesis} & Contradicting evidence \\
\bottomrule
\end{tabular}
\end{table}

\FloatBarrier
\end{block}
\end{frame}

\end{document}
